
\textbf{1.} A partir de la ecuación teórica para el sistema masa-resorte estático:
\begin{equation}
    y = \frac{g}{k}\,m + y_{0}
\end{equation}

\textbf{Responde:} ¿Qué representa la pendiente en esta ecuación? ¿Y la ordenada en el origen?
Si se grafica posición y (m) vs masa m (kg), ¿cómo se obtiene la constante elástica
k a partir de la pendiente?
Supón que en un experimento estático se obtiene una pendiente a = 0,245 m/kg. Calcula
k usando $g = 9,8 m/s^ 2$


\textbf{Respuesta:}
\textbf{¿Qué representa la pendiente en esta ecuación?}
La pendiente corresponde a $\frac{g}{k}$ y representa cómo cambia la posición del
resorte por cada unidad de masa correspondiente. En consecuencia, se puede predecir
el comportamiento del sistema. 

\textbf{¿Y la ordenada en el origen?}
La ordenada en el origen, correspondiente originalmente a b en la ecuación de la recta, 
acá se representa por $y_0$, esto se refiere a la posición inicial del resorte cuando
no hay ninguna masa colgado. 

\textbf{Si se grafica posición y (m) vs masa m (kg), ¿cómo se obtiene la constante 
elástica k a partir de la pendiente?}

Al graficar la masa $m$ en el eje x y la posición $y$ en el eje y, la pendiente
obtenida será $a= \frac{g}{k}$. Al despejar, se obtiene: 
\begin{equation}
 k=\frac{g}{a}  
\end{equation}
 
\textbf{Supón que en un experimento estático se obtiene una pendiente $a = 0,245 m/kg$ 
Calcula $k$ usando $g = 9,8 m/s^ 2$}


\begin{pycode}
g= 9.8
a = 0.245
k = g / a

\end{pycode}

Para calcular k se sustituen los valores en la ecuación 2, donde se obtiene que 
la constante elástica es $k = \py{round(k, 3)}$ N/m.



\textbf{2.} Para el sistema dinámico, el período de oscilación está dado por: 

\begin{equation}
    T = 2 \pi \sqrt{\frac{m+m_e}{k}}
\end{equation}

 · Despeja $T^2$ en función de $m$. Escribe la ecuación lineal resultante (forma $y=ax+b$),
identificando claramente la variable dependiente e independiente, la pendiente y la
ordena en el origen.


\textbf{Respuesta:}
Elevando al cuadrado y despejando se obtiene la forma lineal
\begin{equation}
    T^2 = \frac{4\pi^2}{k}\,m + \frac{4\pi^2 m_e}{k}
\end{equation}

Aquí la variable dependiente es $T^2$, la independiente es $m$, la pendiente es
$a=\dfrac{4\pi^2}{k}$ y la ordenada en el origen es $b=\dfrac{4\pi^2 m_e}{k}$.


 · Si al graficar $T^2$ $vs$ $m$ se obtiene una pendiente $a = 4,0\,\mathrm{s^2/kg}$, 
 ¿cuál es el valor de $k$?

 \begin{pycode}

import math
a = 4.0
k_dos = 4* math.pi**2 / a

 \end{pycode}


\textbf{Respuesta:}

Si $a=\dfrac{4\pi^2}{k}$ entonces; $k=\dfrac{4\pi^2}{a} =
\py{round(k_dos, 2)} N/m$
 
 ·Si además la ordenada en el origen es $b = 0,12 s^ 2$ , ¿cuál es el valor de la masa equivalente
$m_e$ del resorte?

Como ya se tiene a, $b = a \cdot m_e$, despejando: $m_e= \frac{b}{a}$
$m_e = 0.03$ kg


 \textbf{3.} Responde brevemente:
 · ¿Qué significa que el parámetro $R^2$ sea igual a 0.998 en un ajuste lineal?


\textbf{Respuesta:} Significa que el $99.8\%$ de la variación de la 
variable dependiente se explica con una relación lineal con la 
variable independiente, por lo que el modelo lineal es adecuado. Sin
embargo, se debe analizar junto a la representación gráfica de residuos
para prevenir sesgos en el modelo. A su vez, la variación, se refiere
 a qué tanto difieren los valores de la variable dependiente 
respecto a su promedio; si las variables cambian de forma 
linealmente predecible, el valor de $R^2$ se acerca a 1. 

 · ¿Qué tipo de función (lineal, potencial, exponencial)
usarías para graficar $T^2$ vs m según
la teoría?

\textbf{Respuesta:}
De acuerdo a los despejes anteriores, la ecuación para 
esta variable toma la forma de una ecuación de la recta, 
por lo tanto, una función lineal es la más apropiada para 
predecir tendencias con una pendiente de $a = \frac{4\pi^2}{k}$

 · En la práctica anterior se usó papel milimetrado y ajuste visual. ¿Qué ventaja ofrece el
método de mínimos cuadrados con hoja de cálculo?

\textbf{Respuesta:} Este método busca minimizar la suma de los errores 
cuadrados de todos los puntos simultáneamente repecto a la recta, lo que 
obtiene la mejor pendiente e intercepto. Además, elimina el sesgo visual 
y brinda la herramienta de $R^ 2$ para cuantificar la calidad del ajuste.

\textbf{4.} La práctica indica que la masa equivalente del resorte me es igual a un tercio de la
masa real del resorte.
Supón que el resorte utilizado tiene una masa Mresorte = 60 g. Calcula me . Luego, si
en la experiencia dinámica se usó una masa colgante m = 0,200 kg, ¿cuál es la masa
total oscilante que debe usarse en la fórmula teórica del período?

\textbf{Respuesta:} La masa total oscilante es la suma de la masa colgante 
y la masa equivalente; si la masa resorte es de 60g,
 la masa equivalente es de 20g.
 Entonces, para obtener 
la masa total oscilante, se tiene que 
  $m_t = m + m_e = 0.200 + 0.020 = 0.220 \, \text{kg}$

 \textbf{5.} Diseña las tablas de datos que utilizarás en el laboratorio para:
Régimen estático: columnas para masa (g y kg), posición (cm y m).
 Incluye una fila
para $y_0$ sin masa.
Régimen dinámico: columnas para masa (g → kg), 
tiempo de 20 oscilaciones (s),
período T (s), $T ^2$ ($s^2$).

\begin{table}[h]
\centering
\caption{Régimen Estático}
\begin{tabular}{ccccc}
\toprule
$m$ (g) & $m$ (kg) & Posición (cm) & Posición (m) & Elongación $x$ (m) \\
\midrule
0 (sin masa) & 0 & & & 0 \\
& & & & \\
& & & & \\
& & & & \\
& & & & \\
\bottomrule
\end{tabular}
\label{tab:estatico}
\end{table}

\begin{table}[h]
\centering
\caption{Régimen Dinámico}
\begin{tabular}{ccccc}
\toprule
$m$ (g) & $m$ (kg) & $t_{20}$ (s) & $T$ (s) & $T^2$ ($\text{s}^2$) \\
\midrule
& & & & \\
& & & & \\
& & & & \\
& & & & \\
& & & & \\
\bottomrule
\end{tabular}
\label{tab:dinamico}
\end{table}